% Folder 79, batch 09 (archivists' pages 161-180).
% Pass: Opus 5.5 (claude-opus-5-5), 2026-10-03 — first pass, unchecked against the pages by a human.
%
% Notation fixed for the batch: \mathfrak{S}_S the symmetric group of S;
% \mathfrak{P}(S), \mathfrak{P}_2(S), \mathfrak{P}_3(S) the parts (of cardinal 2, 3);
% \mathrm{Et}_A(s) the star of s in the graph A; \mathcal{E} a set of stars.
% The author's own pagination (82, 83, ...) stands on the odd sheets only.
\documentclass[11pt,a4paper]{article}
\input{../preamble/grothendieck}

\folder{79}
\batch{9}
\pages{161}{180}
\foldertitle{2-polyèdres réguliers triangulés, et modules libres de rangs 2 sur les anneaux $\Lambda=\mathbb{Z}/n\mathbb{Z}$ et $\Lambda=\mathbb{Z}$ : notes manuscrites (s.d.).}
\dating{[à partir de 1980]}
\watermark{Édition de démonstration}

\begin{document}

\page{161}\note{p. 82 de l'auteur. La page s'ouvre au milieu d'un argument commencé avant le lot : les numéros (7) et (8) supposent une numérotation antérieure, et $Z$, $\Pi$, $A_0$, $A_i$ y sont déjà fixés.}

\marginal{automorphismes stricts d'un graphe : homothéties (7)}
\marginal{automorphismes stricts d'un multigraphe : homothéties (8)}

\[
G_0 = G_{A_0} = \mathrm{Aut}_Z(S,A_0,Z)
= \Bigl\{ g\in\mathfrak{S}_S \Bigm|
\begin{array}{l}
g \text{ autom. de } (S,A_0) \text{ i.e. } g(A_0)=A_0\\
g \text{ commute à } Z
\end{array}\Bigr\}
\]
\note{au-dessus de $G_{A_0}$, sous une accolade : « automorphismes du gr. \uncertain{à} homothéties ».}

\[
G = \mathrm{Aut}_Z(S,\Pi)
= \Bigl\{ g\in\mathfrak{S}_S \Bigm|
\begin{array}{l}
g \text{ autom. de } (S,\Pi) \text{ i.e. } g(\Pi)=\Pi\\
g \text{ commute à } Z
\end{array}\Bigr\}
\]
i.e. l'opération de $g_\Pi$ sur $\Pi$ se fait par translation par un él. de $Z$
\[
= \Bigl\{ g\in\mathfrak{S}_S \Bigm|
\begin{array}{l}
g \text{ commute à } Z\\
g(A_0)\in\Pi
\end{array}\Bigr\}.
\]
\note{sous l'accolade de $G$ : « \ill{} des multigraphes : homothéties ».}

$G_{A_i}$ indépendant de $i\in\Pi$.

\struck{\ill{}} \ill{}

\begin{tikzcd}
1 \arrow[r] & G_0 \arrow[r, hook] & G \arrow[r] & Z/\mathfrak{z} \arrow[r] & 1 \\
 & & Z \arrow[u] \arrow[ur, "\lambda\mapsto\lambda^2"'] & &
\end{tikzcd}

\note{le quotient est écrit avec une lettre en forme de 3 bouclé, transcrite ici $\mathfrak{z}$ (cf. la ligne suivante) ; au-dessus, un symbole biffé illisible.}

\[
\mathfrak{z} = \mathrm{Ker}\bigl(\lambda\mapsto T_\lambda : Z\to\mathfrak{S}_\Pi\bigr)
= \bigl\{\lambda\in Z \bigm| \mathrm{Et}_{A_i}(s)=\mathrm{Et}_{A_i}(\lambda s)\ \ \forall s\in S,\ \forall i\in\Pi\bigr\}
\]
\[
\mathfrak{z}=\{1\} \text{ si tout sommet est déterminé par son étoile …}
\]

\struck{NB (12) Soit}
\[
\mathcal{E}_{A_0} = \bigl\{E\in\mathfrak{P}(S) \bigm| \exists s\in S\ ;\ E=\mathrm{Et}_{A_0}(s)\bigr\}
\]
$\mathcal{E}_{A_i}$ indépendant de $i\in\Pi$, soit $\mathcal{E}=\mathcal{E}_\Pi$.

\struck{étoiles} \uncertain{conservant} $\mathcal{E}\subset\mathfrak{P}(S)$ (des étoiles pour les graphes $A_i$, mais le \uncertain{discours} \uncertain{structurant} des « centres » \struck{\ill{}} (ou les centres) d'une \uncertain{étoile} $E\subset S$ vérifiant … \ill{} la structure de \ill{} était fixée). Pour tout $g\in G_0$ on a $g(\mathcal{E})=\mathcal{E}$ i.e. $g$ transforme étoile en étoile. C'est \uncertain{en} \uncertain{particulier} pour les $\lambda\in Z$.

\marginal{NB Il faudrait quand même expliciter ici des cas où la structure $(S,\Pi)$ implique déjà la \uncertain{connaissance} \uncertain{(9)} de $Z$, et \uncertain{aussi} \struck{\ill{}} des cas où \ill{} \ill{} \uncertain{alors} $Z$ \ill{} $(S,\Pi)$ \ill{} \ill{} déjà à \ill{} trivial \ill{} \uncertain{congruentielle} \ill{} induit l'autre (10) \ill{} (c'est \ill{} cas \ill{} les \ill{} : il suffit \ill{} : $Z$ \ill{}) \ill{} $(S,A_0)$ \ill{} tout \ill{} de $G_0$ \ill{} (11) C'est le cas \ill{} dont il \ill{} \ill{} dans \ill{} \ill{} prépolyédral \ill{} \uncertain{NB} \uncertain{retrouve} $(S,F_0)$ \ill{} $\Pi$ \ill{} \ill{} 2-polyèdre str. régulier \ill{} \ill{} $(G_0)$ \ill{} (12) \ill{} $(S,F)$, \ill{} $(S,\Pi)$, \ill{} centre de $G$. En \ill{}, \ill{} (13) \ill{}}
\note{ce long ajout est écrit en diagonale dans la moitié gauche de la feuille, par-dessus le reste ; on n'en lit que des fragments. Il porte les numéros (9) à (13).}

\page{162}
On dit qu'un graphe $(S,A_0)$ est \emph{régulier} si $G_0$ (le groupe de ses automorphismes stricts) est transitif sur l'ens. $\vec{A}_0$ de ses arêtes \add{et si $(S,A_0)$ n'a pas de sommet isolé et est connexe} \uncertain{orientées}. On dit qu'un multigraphe \struck{\ill{}} à homothéties est régulier, si les graphes constituants le \struck{sont} \add{sont \uncertain{réguliers}}.

Si $\mathcal{G}(S,A_0,Z)$ est \uncertain{régulier}, $G_0$ étant transitif sur $S$, $Z$ opère \emph{librement} sur $S$. En \uncertain{tous} cas, j'appelle \emph{\uncertain{voyage}} des graphes \uncertain{à} homothéties, toute \ill{} de $Z$ dans $S$.

\struck{\ill{}} Si $(S,\Pi,Z)$ \struck{est un multigraphe \ill{}} est régulier, alors \struck{pour $i,j\in\Pi$}
\[
A_i\cap A_j=\emptyset \quad\text{si } i,j\in\Pi,\ i\neq j \tag{14}
\]
car si $a\in A_i\cap A_j$, alors $A_i$ est déterminé \uncertain{par} $a$ \ill{} l'orbite de $a\in\mathfrak{P}_2(S)$ sous $G_0$.

On \uncertain{trouve} \ill{} une \uncertain{application} surjective
\[
A \overset{\text{déf}}{=} \bigcup_{i\in\Pi} A_i \xrightarrow{\ \tau_A\ } \Pi \tag{15}
\]
\uncertain{définie} par une partition
\[
A=\coprod_{i\in\Pi} A_i .
\]

\marginal{La structure de multigraphe $(S,\Pi,Z)$ (15) équivaut à \uncertain{une} structure $(S,A,R,Z)$, \ill{} $A\subset\mathfrak{P}_2(S)$, $R$ \uncertain{rel. d'équivalence} \ill{} (16) \ill{} (plus \ill{} des \ill{} $(S,A_i,Z)$ \ill{}) \ill{} (17) \ill{}}
\note{ajout en diagonale dans la marge gauche, peu lisible ; il porte les numéros (15) à (17).}

\struck{Soit \ill{} On dit qu'un} Revenons à un graphe \uncertain{régulier} $(S,A_0)$, et posons
\[
F_0=\bigl\{f\in\mathfrak{P}_3(S) \bigm| \forall a\in\mathfrak{P}_2(f),\ a\in A_0\bigr\}
\]
\[
A_0 \equiv \bigl\{a\in\mathfrak{P}_2(S) \bigm| \exists f\in F_0,\ f\supset a\bigr\}
\]
d'où une structure de \struck{triangulation} « \emph{prépolyèdre} » \uncertain{du} graphe, dits \struck{\ill{}} \ill{} si
\[
P_0=(S,A_0,F_0,\text{ relations d'incidence évidentes}) \tag{18}
\]

\page{163}\note{p. 83 de l'auteur.}
est bel et bien un 2-polyèdre — i.e. s'il satisfait la condition \struck{de connexité locale \ill{}} \struck{\ill{} que \ill{} et}
\[
\forall a\in A_0,\ a=\{s,t\},\ \text{il existe exactement deux } u\in S \text{ telles que } \{s,u\},\{t,u\}\in A_0 \tag{19}
\]
plus la condition de connexité des contours locaux en tout $s\in S$. S'il en est ainsi, $G_0$ opère bien \uncertain{sur} \uncertain{la} structure polyédrale. Si le graphe \uncertain{à} homothéties est régulier, \struck{\ill{}} il en est de même \uncertain{du} 2-polyèdre associé \ill{}; \uncertain{plus} \ill{} si le graphe à homothéties est prépolyédral, et si $G_0$ \ill{} simplement transitif) \uncertain{on} \ill{} des 2-polyèdres associés \ill{} les multigraphes \uncertain{réguliers} \ill{} \ill{} str. \emph{réguliers} : alors $G_0$ est simplement transitif sur l'ens. des repères de $(S,F)$ donc une description \ill{} en termes de théorie des groupes \ill{} en termes \ill{} théorie combinatoire \ill{} \uncertain{voyage} \ill{}, qui est bien comprise. Je \struck{\ill{}} fixe un repère \struck{\ill{}}
\[
(s_0,t_0,u_0)\in\mathrm{Rep}(P_0) \tag{20}
\]
pour \uncertain{expliciter} la situation \uncertain{entièrement}. \struck{D\ill{}} On veut

\marginal{Dans le cas \ill{} : $F=\coprod_{i\in\Pi}F_i$ \ill{} structure de multigraphe \ill{} $(S,F)$, $F\subset\mathfrak{P}_3(S)$, $Z\subset$ \ill{} transitif \ill{} \ill{} j'ai envie de \ill{}}
\note{ajout en diagonale dans le coin inférieur gauche.}

\page{164}
tout décrire en termes de théorie des groupes, \struck{\ill{}} \add{en termes d'une structure} de groupes \ill{} l'intention (10). Ici on trouve des générateurs canoniques de $G_0$
\[
\left\{
\begin{array}{l}
\tau_0,\tau_1,\tau_\infty\in G_0 \ \text{satisfaisant}\\
\rho_1^2=\rho_\infty^3=1 \quad\text{où}\quad \rho_1=\tau_0\tau_\infty,\ \rho_\infty=\tau_1\tau_0,\\
\phantom{\rho_1^2=\rho_\infty^3=1 \quad\text{où}\quad} \rho_0=\tau_\infty\tau_1
\end{array}\right. \tag{21}
\]
si $n$ est l'ordre de $\rho_0$ (il sera fini dans les cas \uncertain{qui} nous \uncertain{intéressent}) on aura aussi
\[
\rho_0^{\,n}=1 , \tag{22}
\]
$n$ sera aussi l'ordre des sommets du graphe régulier $\Gamma_0=(S,A_0)$, i.e. des 2-polyèdres réguliers $P_0$. Il faut \uncertain{maintenant} prendre $S$ comme espace homogène sous $G$, via l'origine $s_0\in S$, donc il faut \uncertain{déterminer} le stabilisateur
\[
G_{s_0}=\{g\in G\mid gs_0=s_0\}=B' \tag{23}
\]
dont l'intersection avec $G_0$ sera donnée par
\[
B'\cap G_0 = \text{sous-groupe engendré par } \tau_1,\tau_\infty = \tilde{U}\ \ (\text{isomorphe à } \mathbf{D}_n) \tag{24}
\]
on a donc
\[
\begin{array}{rcl}
G/B' & \xrightarrow{\ \sim\ } & S\\
g.B' & \longmapsto & g.s_0
\end{array}
\qquad\text{isom. d'espaces homogènes sous } G \tag{25}
\]
\note{les numéros (23), (24), (25) en marge sont surchargés ; ils se lisent d'après la suite (21), (22).}

\drawing{un repère : les sommets $s_0$, $t_0$, $u_0$, $u'_0$ ; le triangle $s_0t_0u_0$ est découpé en faces étiquetées $r$ (hachurée), $\tau_1r$, $\tau_0r$, et sous l'arête $s_0t_0$ une face $\tau_\infty r$ en pointillé.}

\page{165}\note{p. 84 de l'auteur.}
\struck{Notons que pour un espace homogène}

Traduisons en termes de théorie des groupes, de façon tout à fait générale, la donnée suivante :
\begin{enumerate}
\item[(26)] a) un espace homogène $S$ sous un groupe $G_0$, avec « origine » fixée $s_0\in S$,

b) surgroupe $G\supset G_0$ opérant sur $S$ (en prolongeant l'action précédente), avec $G_0$ distingué dans $G$ (soit $Z'=G/G_0$),

c) sous-groupe central $Z$ de $G$,
\end{enumerate}
\uncertain{de la} \uncertain{manière} \uncertain{la plus} possible, de faire une « traduction » qui soit « produite par le groupe $G_0$ », et de groupes qui se \uncertain{décrivent} aisément en termes de $G_0$, \struck{\ill{}} ses sous-groupes etc.

\marginal{NB On a \uncertain{pris} (26), dans cette \uncertain{présentation} \ill{} si les groupes $G_0$, $G$ opérant fidèlement sur $S$ \ill{} \ill{} \ill{}}

La donnée a) revient à la donnée du couple
\[
G_0\supset\tilde{U} \tag{27}
\]
d'un groupe $G_0$, et d'un sous-groupe $\tilde{U}$ $(=G_{0s_0})$, qui permet de reconstruire $(S,s_0)$ comme
\[
(S,s_0)\simeq(G_0/\tilde{U},\ \dot{1}_G) \tag{28}
\]
\note{sous $\dot 1_G$ : « classe de $1_G$ dans $G_0/\tilde U$ ».}

La donnée b) s'exprime, tout d'abord, comme celle d'un diagramme de groupes

\begin{tikzcd}
G_0 \arrow[r, hook, "\text{dist.}"] & G \\
\tilde{U} \arrow[u, hook] \arrow[r, hook] & B' \arrow[u, hook]
\end{tikzcd}

\[
(29) \qquad (\mathrm{NB}\ \ B'=G_{s_0})
\]
\note{le dessin met $G_0$ et $G$ en haut, $\tilde U$ et $B'$ en bas, les flèches verticales descendant ; elles sont rendues ici comme inclusions montantes.}

« cartésien et cocartésien », cartésien signifiant
\[
\tilde{U}=G_0\cap B' \quad (\text{donc } \tilde{U} \text{ dist. dans } B') \tag{30}
\]
« cocartésien » s'exprime par une des deux conditions équivalentes
\[
\left\{
\begin{array}{ll}
G_0/\tilde{U} \xrightarrow[\text{can}]{\ \sim\ } G/B' & \text{qui équivaut à}\\[2pt]
B'/\tilde{U} \xrightarrow{\ \sim\ } G/G_0 & (\text{isom. de groupes}).
\end{array}\right. \tag{31}
\]
Le groupe $G$, extension d'un groupe

\page{166}
\[
Z'=G/G_0 \tag{32}
\]
par $G_0$, est \uncertain{plutôt} \uncertain{mal} « éloigné » \uncertain{directement} de $G_0$, on préfère l'exprimer en termes des groupes $B'$, extension d'une \uncertain{même} $Z'$ par le \add{sous-}groupe \struck{$M$} $\tilde{U}$ de $G$ (en pratique, $Z'\simeq(\mathbb{Z}/n\mathbb{Z})^*$ ou $(\mathbb{Z}/n\mathbb{Z})^*/\pm1$, $\tilde{U}\simeq\mathbb{Z}/n\mathbb{Z}$ ou $\mathbf{D}_n$, sont deux groupes bien compris et « petits », leur extension sera aussi bien \uncertain{gérable}). En outre des \uncertain{tapis} bien connus, pour reconstituer l'extension $G$ de $Z'$ par $G_0$ : l'idée de l'extension $B'$ de $Z'$ par $\tilde{U}$, par \uncertain{voie} \uncertain{naturelle} de \struck{\ill{}} \add{extension} des groupes \uncertain{régis} $\tilde{U}\hookrightarrow G_0$, il faut de plus se donner l'opération
\[
B'\longrightarrow \mathrm{Aut}_{\mathrm{gr}}(G_0) \tag{33}
\]
prolongeant l'opération \uncertain{tautologique} des sous-groupes $\tilde{U}$ de $B'$ sur $G_0\supset\tilde{U}$, par \uncertain{automorphismes} intérieurs.

Ainsi, les données (a)(b) de (26) équivalent aux données
\begin{enumerate}
\item[(34)] a) une extension $B'$ d'un groupe $Z'$ par un groupe $\tilde{U}$
\[
1\to\tilde{U}\to B'\to Z'\to 1
\]
b) un hom. injectif de groupes
\[
\tilde{U}\hookrightarrow G_0
\]
c) une opération (33) de $B'$ sur $G_0$, prolongeant l'opération de $\tilde{U}$ déduite de b).
\end{enumerate}

Pour interpréter la donnée c) de (26), on \uncertain{commence} : notons \struck{que le commutant} \add{l'isomorphisme de} groupes

\page{167}\note{p. 85 de l'auteur.}
\[
\text{Commutant de } G_0 \text{ dans } \mathfrak{S}_S \ \simeq\ \mathrm{Norm}_{G_0}(\tilde{U})/\tilde{U} \tag{35}
\]
donc la donnée d'un sous-groupe $Z$ du commutant revient à la donnée d'un sous-groupe $\tilde{B}''$ de $G_0$, \uncertain{satisfaisant}
\[
\tilde{U}\subset\tilde{B}''\subset\mathrm{Norm}_{G_0}(\tilde{U}),\quad\text{avec } Z= \tag{36}
\]
i.e. $\tilde{B}''$ un sous-groupe de $G_0$, contenant $\tilde{U}$ comme sous-groupe \emph{distingué} \uncertain{de sorte que} \ill{} \uncertain{groupe} \uncertain{quotient}
\[
Z=\tilde{B}''/\tilde{U}
\]
\note{la fin de (36), « $Z=\tilde B''/\tilde U$ », est écrite à droite dans un cadre relié à la ligne par un trait.}
en termes de (26), on a
\[
\tilde{B}''=\{g\in G_0\mid gs_0\in Z.s_0\}. \tag{37}
\]
La condition que $Z$ commute non seulement à $G_0$, mais aussi à l'action de $G$ toute entière, ce qui revient au même de $B'$ (puisque $G=B'.G_0$ produit des 2 sous-gr. de $G$) s'explicite aisément par l'une des conditions équivalentes
\begin{enumerate}
\item[(38)] a) $\tilde{B}''$ normalise le sous-groupe $B'$ de $G$, \uncertain{maximum} :

b) $B'$ normalise $\tilde{B}''$ (pour l'action (33)) et l'action dans le quotient $\tilde{B}''/\tilde{U}$ est triviale.
\end{enumerate}
Ici c'est bien sûr la forme 38 b) qui sera la plus commode pour notre propos, \struck{$Z$} s'exprimant directement en termes des données (34) et (36).

Nous pouvons introduire alors le \add{sous-}groupe $B$ de $G$ engendré par $\tilde{B}''$, $B'$, qui s'insère dans le diagramme

\begin{tikzcd}
G_0 \arrow[r, no head, "{Z'}"] \arrow[d, no head] & G \arrow[d, no head] \\
\tilde{B}'' \arrow[r, no head, "{Z'}"] \arrow[d, no head, "Z"'] & B \arrow[d, no head, "Z"] \\
\tilde{U} \arrow[r, no head, "{Z'}"] & B'
\end{tikzcd}

\[
(39)
\]
où on a marqué au-dessus des flèches d'inclusion, d'

\page{168}
\uncertain{un} sous-groupe distingué dans un \uncertain{groupe} \uncertain{ambiant}, la valeur du groupe quotient. Les deux carrés de ce diagramme sont cartésiens et cocartésiens — en particulier, $B$ \uncertain{n'est} \uncertain{pas} \uncertain{un} sous-groupe \struck{\ill{}} \uncertain{inv.} $\tilde{B}''$, se déduit de $(B'\supset\tilde{U})$ et de l'action de $B'$ sur $\tilde{B}''$ (induite par (33)), de même que $(G\supset G_0)$ se déduit de $(B'\supset\tilde{U})$ et de l'action (33) de $B'$ sur $G_0$. Le groupe $\tilde{B}''$ \struck{\ill{}} \uncertain{sera}, dans les cas qui nous occupent, \add{(en plus d'être \uncertain{plongé} dans $G_0$)} « petits » et bien \uncertain{figurables}, au même titre que $B'$ — quel \ill{} il \uncertain{vraisemblablement} \ill{} \uncertain{fin}, et il en sera de même de $B$.

Nous avons exprimé par l'introduction \add{dans (26 c))} de la donnée \ill{} sous-groupe $Z$ de $\mathfrak{S}_S$ commutant à $G$, mais il vaut mieux : \uncertain{exprimer} que $Z$ est donné en tant que sous-groupe central de $G$. Cela s'explicite par la donnée d'un homomorphisme \struck{\ill{}}
\[
Z\longrightarrow G \tag{40}
\]
tel que les deux actions de $Z$ sur $S$ — celle donnée \uncertain{au départ}, et celle déduite de (40), soient égales. On voit alors tout de suite que cela signifie que (40) provient d'un \emph{scindage} (de l'extension $B$ de $Z$ par $B'$, \struck{\ill{}}
\[
\bullet\longrightarrow B'\longrightarrow B\longrightarrow Z\longrightarrow\bullet \tag{41}
\]
définissant (\uncertain{donc}) un isomorphisme $B\simeq Z\times B'$, avec en plus $Z\subset B$ opérant \emph{trivialement} sur $G_0$ (donc sur $G$).
\note{une flèche courbe revient de $Z$ vers $B$ : c'est le scindage.}

\page{169}\note{p. 86 de l'auteur.}
Nous allons reformuler les données (34)(36) et (41) (le scindage $\xi$ de l'extension $B$ de $Z$ par $B'$) \struck{en \ill{}} \add{comme une \ill{}} \add{des} variantes \uncertain{précises} \struck{de diagrammes} des données suivantes :

(42) a) Un diagramme \add{commutatif} de groupes, cartésien et cocartésien, avec des inclusions de groupes qui se distinguent,

\begin{tikzcd}
\tilde{B}'' \arrow[r, hook, "{Z'}"] & B \\
\tilde{U} \arrow[u, hook, "Z"] \arrow[r, hook, "{Z'}"] & B' \arrow[u, hook, "Z"']
\end{tikzcd}

\note{dans le manuscrit la flèche de gauche descend de $\tilde B''$ vers $\tilde U$ ; celle de droite monte de $B'$ vers $B$. Les deux sont rendues ici montantes, comme inclusions.}

b) Un scindage de l'extension $B$ de $Z$ par $B'$ ; \uncertain{avec} un groupe-section $Z\subset B$ central dans $B$, de sorte que l'on a un iso. d'ext.
\[
B\simeq B'\times Z,\qquad Z\simeq\tilde{B}''/\tilde{U} \ \text{commutatif}
\]
\note{devant le second isomorphisme, un symbole biffé.}

c) Une \struck{opération \ill{}} \add{injection} de groupes
\[
\tilde{B}''\hookrightarrow G_0
\]
d) Une opération de $B$ sur $G_0$,
\[
B\longrightarrow\mathrm{Aut}_{\mathrm{gr}}(G_0)
\]
prolongeant l'opération tautologique de $\tilde{B}''$ sur $G_0$ déduite de c), et \emph{triviale} sur le sous-groupe $Z$ de $B$. (Opération qui permet de reconstituer le groupe $G$ et le diagramme (39).)

Nous allons encore modifier un peu cette traduction, en écrivant le diagramme (42 a) \add{+ donnée b)}, sous la forme

\begin{tikzcd}
\tilde{B}'' \arrow[r, hook, "t"'] & B'\times Z \\
\tilde{U} \arrow[u, hook, "Z"] \arrow[r, hook, "{Z'}"] & B'=B'\times\{1\} \arrow[u, hook]
\end{tikzcd}

\[
(43)\qquad Z=\tilde{B}''/\tilde{U}
\]
\note{au-dessus de la flèche $t$, un symbole noirci ; « $B'\times\{1\}$ » est une lecture incertaine du second facteur.}

\page{170}\note{p. 87 de l'auteur.}
qui se reconstitue à partir du diagramme de groupes partiel

\begin{tikzcd}
\tilde{B}'' & \\
\tilde{U} \arrow[u, hook, "Z"] \arrow[r, hook, "{Z'}"] & B'
\end{tikzcd}

\[
(44)
\]
sauf que pour en tirer l'hom. horizontal $i$ de (43), dont la composante suivant $Z$ doit être l'hom. canonique dans le groupe quotient, il faut se donner la première composante, savoir un homomorphisme de groupes
\[
\tilde{B}''\longrightarrow B'
\]
la commutativité de (43) signifiant que ce dernier rend commutatif le triangle

\begin{tikzcd}
\tilde{B}'' \arrow[dr] & \\
\tilde{U} \arrow[u, hook, "Z"] \arrow[r, hook, "{Z'}"] & B'
\end{tikzcd}

\[
(45)
\]
de telle sorte qu'elles \add{donc} donnent naissance à un hom. d'extensions de groupes

\begin{tikzcd}
1 \arrow[r] & \tilde{U} \arrow[r] \arrow[d, no head, "\shortparallel" description] & \tilde{B}'' \arrow[r] \arrow[d, "\psi"] & Z \arrow[r] \arrow[d] & 1 \\
1 \arrow[r] & \tilde{U} \arrow[r] & B' \arrow[r] & {Z'} \arrow[r] & 1
\end{tikzcd}

\[
(46)
\]
Ainsi, les données (42 a) b)) se réduisent à la donnée d'un hom. de suites exactes (46), ou encore d'un diagramme d'hom. de groupes
\[
(47)\qquad
\begin{array}{ccc}
 & & Z\\
 & & \downarrow\\
B' & \xrightarrow[\text{épi}]{} & Z'
\end{array}
\]
\note{la suite de (47) passe à la page suivante.}

\page{171}\note{p. 88 de l'auteur.}
dont on déduit \struck{$\tilde{U}\simeq\mathrm{Ker}(B'\to Z')$,} \ill{} et $\tilde{B}''$ comme
\[
\tilde{B}''\simeq B'\times_{Z'}Z ,\qquad \tilde{U}=\mathrm{Ker}(B'\to Z'). \tag{48}
\]
\note{sous le produit fibré, une accolade : « \uncertain{produit} \uncertain{fibré} de l'ext. $B'$ de $Z'$ par $\tilde U$, par l'hom. $Z\to Z'$ ».}

\struck{Ceci dit,} Il nous semble que c'est sous la forme (46), \uncertain{et} \uncertain{non} \uncertain{autrement}, que la donnée $\tilde{B}''$ figure \uncertain{déjà} \add{\uncertain{cependant}} nécessairement, et le plus naturelle, puisque c'est \uncertain{sous} $\tilde{B}''$ que interviennent directement dans les données (42 c)). Pour expliciter enfin la donnée (41 d), en écrivant $B=B'\times Z$, et en notant que l'opération de $B$ sur $G_0$ doit être triviale sur $Z$, on voit que cette donnée se réduit à celle d'une op. de $B'$ sur $G_0$ \struck{i.e. \ill{} (33)} pour écrire que \struck{cette} l'opération \add{de $B$ sur $G$} se réduit sur $\tilde{B}''$ à l'opération tautologique déduite du plongement $\tilde{B}''\subset G_0$, on note que l'opération \uncertain{composée}
\[
\tilde{B}''\longrightarrow B'\longrightarrow\mathrm{Aut}_{\mathrm{gr}}(G_0) \tag{49}
\]
qui doit donc être l'opération tautologique.

En résumé, on trouve que les données (26) s'explicitent en termes de théorie des

\page{172}
groupes par les données suivantes

(50) a) Un diagramme de groupes
\[
\tilde{U}\subset\tilde{B}''\subset G_0,\quad\text{avec } \tilde{U} \text{ distingué dans } \tilde{B}''
\]
\struck{\ill{}} d'où un groupe quotient
\[
Z=\tilde{B}''/\tilde{U}\qquad(\tilde{B}'' \text{ ext. de } Z \text{ par } \tilde{U})
\]
b) Un \struck{hom.} homom. de groupes $Z\xrightarrow{\ \varepsilon\ }Z'$ et une extension \add{$B'$} de $Z'$ par le $\tilde{U}$ de a), de telle façon que $\tilde{B}''$ soit déduit de cette extension par l'hom. $Z\to Z'$, i.e. qu'on ait un hom. de suites exactes (46).

c) Une opération de $B'$ sur $G_0$, de telle façon que l'opération composée
\[
\tilde{B}''\xrightarrow{\ \psi\ }B'\longrightarrow\mathrm{Aut}_{\mathrm{gr}}(G_0)
\]
soit l'opération tautologique déduite du plongement a) de $\tilde{B}''$ dans $G_0$.

\marginal{oubli : la condition que pour $g\in\tilde{B}''$, $\theta(g)\,|\,\tilde{U}=\mathrm{int}(g)\,|\,\tilde{U}$, et que $\theta(g)$ \uncertain{normalise} $\tilde{U}$ et centralise $\tilde{B}''$ mod $\tilde{U}$}
\note{ajout en diagonale à gauche de c), relié à l'accolade de (50).}

En termes des données (26), on a :
\[
(51)\quad\left\{
\begin{array}{l}
\tilde{U}=G_{0s_0}=\{g\in G_0\mid g.s_0=s_0\}\\
\tilde{B}''=\{g\in G_0\mid gs_0\in Z.s\}\\
Z'=G/G_0,\quad Z\to Z' \text{ le composé } Z\xrightarrow[\text{incl.}]{}G\xrightarrow{\text{can}}G/G_0=Z'\\
B'=G_{s_0}=\{g\in G\mid gs_0=s_0\} \text{ de sorte que}\\
\phantom{B'=}\ \tilde{U}=B'\cap G_0,\quad B'\to Z' \text{ composé } B'\xrightarrow[\text{incl}]{}G\xrightarrow{\text{can}}G/G_0=Z'\\
\text{op. de } B' \text{ sur } G_0 \text{ induite par celle de } G \text{ sur le sous-groupe distingué } G_0
\end{array}\right.
\]
\note{dans (51), « $\{g\in G_0\mid gs_0\in Z.s\}$ » est écrit tel quel ($s$ pour $s_0$). Devant $\{g\in G_0\mid g.s_0=s_0\}$ et devant $G_{s_0}$, des symboles noircis ; sous $B'=G_{s_0}$ une ligne commençant par un « $Z=$ » biffé.}

En termes des données (50), on récupère (26) par

(52) a) $G$ ext. de $Z'$ par $G_0$ déduite de \add{b)} l'ext. $B'$ de $Z'$ par $\tilde{U}$, le plongement $\tilde{U}\subset G_0$ et l'opération \add{c)} de $B'$ sur $G_0$.

b) $S\simeq G_0/\tilde{U}\simeq G/B'$

c) $Z=\mathrm{Im}\bigl(\tilde{B}''\xrightarrow{\ \varphi\ }G\bigr)$, \quad $\varphi(g)=g.\psi(g)^{-1}=\psi(g)^{-1}.g$

\marginal{produit dans $G$ qui contient $\tilde{B}''$, $B'$ ; $\psi:\tilde{B}''\to B'$ défini dans (50 c) !}
\marginal{NB L'hom. $\varphi:\tilde{B}''\to G$ se factorise par le quotient $\tilde{B}''/\tilde{U}\simeq Z$}

\page{173}\note{p. 89 de l'auteur.}
\emph{NB} On a un homom. unique (via les données (50))
\[
\varphi_Z : Z=\tilde{B}''/\tilde{U}\longrightarrow G \tag{53}
\]
\struck{défini} défini par factorisation : partir de l'application
\[
\varphi=\varphi_{\tilde{B}''} : \tilde{B}''\longrightarrow G,\qquad \varphi(g)=g\,\psi(g)^{-1}=\psi(g)^{-1}g \tag{54}
\]
\note{la seconde égalité est soulignée d'un crochet.}
l'égalité des deux expressions provient de la relation
\[
\psi(g)\,h\,\psi(g)^{-1}=g\,h\,g^{-1}\quad\text{pour } g\in\tilde{B}'',\ h\in G_0 \tag{55}
\]
(cf. (50 c)). La chose essentielle à vérifier est la multiplicativité de $\varphi$ définie par (54), \uncertain{étant} entendu qu'elle est triviale sur $\tilde{U}$ puisque pour $g\in\tilde{U}$, on a $\psi(g)=g$ d'où $\varphi(g)=1$. \struck{Comme} Il faut vérifier
\[
g\,\psi(g)^{-1}\,g'\,\psi(g')^{-1}\overset{!}{=}gg'\,\psi(gg')^{-1}\qquad g,g'\in B'
\]
\note{sous $g'\psi(g')^{-1}$, une flèche marquée « commutent » le fait passer à gauche.}
i.e.
\[
\underbrace{\psi(g)^{-1}\psi(g')^{-1}}_{\psi(g'g)^{-1}}\overset{?}{=}g'\,\psi(gg')^{-1}g'^{-1}
\]
Comme $Z=\tilde{B}''/\tilde{U}$ est commutatif, on a
\[
g'g=gg'u\quad (\text{i.e. } u=g'^{-1}g^{-1}g'g \text{ soit } u^{-1}=g^{-1}g'^{-1}gg')
\]
\struck{d'où}
\[
\psi(g'g)^{-1}=\psi(gg'u)^{-1}=\psi(u^{-1})\,\psi(gg')^{-1},
\]
et la relation à vérifier s'écrit
\[
\underbrace{\psi(u^{-1})}_{=\,u^{-1}=\,g^{-1}g'^{-1}gg'}\overset{?}{=}g'\,\underbrace{\psi(gg')^{-1}g'^{-1}\psi(gg')}_{(gg')^{-1}g'^{-1}(gg')\ \text{(55)}}
\]
soit
\[
g^{-1}g'^{-1}gg'=g'g'^{-1}g^{-1}g'^{-1}gg' \qquad \text{ok.}
\]
\note{la vérification est faite pour $g,g'\in B'$ selon l'énoncé, mais porte manifestement sur $g,g'\in\tilde B''$.}

\page{174}
Avec commutativité dans

\begin{tikzcd}
Z \arrow[r, "\varphi_Z"] \arrow[dr, "(50\,\mathrm{b})"'] & G \arrow[d, "(52\,\mathrm{a})"] \\
 & {Z'}
\end{tikzcd}

\note{sous la flèche $\varphi_Z$, le renvoi « (53) » ; devant « (52 a) » un caractère noirci.}

Nous allons expliciter maintenant les particularités des données (50), dans le cas qui nous occupe.

(56) 1) Le groupe $G_0$ a des générateurs $\tau_0,\tau_1,\tau_\infty$ satisfaisant les conditions (21)

2) $\tilde{U}=$ sous-groupe engendré par $\tau_1,\tau_\infty$, isomorphe à $\mathbf{D}_n$ (non canoniquement)

3) Opération fidèle de $G$ sur $S$, ce qui s'exprime par la fidélité de l'opération de $B'$ \add{($\simeq G_{s_0}$)} sur $S\simeq G_0/\tilde{U}$, opération déduite de l'opération de $B'$ sur $G_0$, stabilisant $\tilde{U}$ et passant donc au quotient

4) $Z'=Z$ et $Z\to Z'=Z$ est l'hom $\lambda\mapsto\lambda^2$, en d'autres termes la donnée b) de (50) est une façon d'exprimer l'extension $\tilde{B}''$ de $Z$ par $\tilde{U}$ comme déduite d'une \add{\uncertain{par image inverse}} extension \add{$B'$} de $Z$ par $\tilde{U}$, via l'hom
\[
Z\longrightarrow Z^2
\]
(« carré » de l'ext. $B'$).

\marginal{donc $B'$ opère fidèlement sur $G_0$, et peut donc s'interpréter comme un sous-groupe de $\mathrm{Aut}_{\mathrm{gr}}(G_0\supset\tilde{B}''\supset\tilde{U})$ contenant l'image de $\tilde{B}''$ dans ce groupe d'automorphismes (ce qui explicite en même temps $\psi:\tilde{B}''\to B'$)}

Donc on peut expliciter les données ainsi
\note{la page s'arrête là ; la suite est à la page suivante.}

\page{175}\note{p. 90 de l'auteur.}
(57) a) Un groupe quotient $G_0$ du groupe cartographique triangulé
\[
\mathfrak{C}_{\infty,3}\simeq\bigl\{\tau_0,\tau_1,\tau_\infty \bigm| \tau_0^2=\tau_1^2=\tau_\infty^2=(\underbrace{\tau_1\tau_\infty}_{\rho_\infty})^3=(\underbrace{\tau_0\tau_\infty}_{\rho_1})^2\bigr\}
\]
b) Un sous-groupe $\tilde{B}''$ \struck{du normalisateur} \add{de $G_0$} contenant le sous-groupe $\tilde{U}$ engendré par $\tau_1,\tau_\infty$ (on a donc, par $\rho_0=\tau_1\tau_\infty$, et $\tau_1$ ou $\tau_\infty$, $\tilde{U}\simeq\mathbf{D}_n$ pour $n=$ l'ordre de $\rho_0$ dans $G_0$), normalisant $\tilde{U}$

c) Un sous-groupe $B'$ de $\mathrm{Aut}_{\mathrm{gr}}(G_0\supset\tilde{B}''\supset\tilde{U})$, \add{opérant trivialement sur $Z=\tilde{B}''/\tilde{U}$,} contenant l'image de $\tilde{B}''$ dans ce groupe par $g\mapsto\mathrm{int}_{G_0}(g)$, d'où un hom
\[
\psi:\tilde{B}''\longrightarrow B'
\]
[On suppose que $G_0$ opère fidèlement sur $G_0/\tilde{U}$ (i.e. l'intersection des conjugués de $\tilde{U}$ dans $G_0$ est $\{1\}$) et que $B'$ opère fidèlement sur $G_0/\tilde{U}$ par transport de structure — les premières conditions indiquent déjà que $\psi$ induit un plongement $\tilde{U}\xrightarrow{\psi|\tilde{U}}B'$, on suppose $\tilde{U}$ \emph{distingué} dans $B'$, on a donc un quotient et un hom déduit par passage au quotient
\[
Z\overset{\text{déf}}{=}\tilde{B}''/\tilde{U}\xrightarrow{\ \kappa\ }B'/\tilde{U}\overset{\text{déf}}{=}Z'.\ ]
\]
d) Un isomorphisme
\[
Z'\xrightarrow[\sim]{\ \bar{\alpha}\ }Z
\]
rendant commutatif le diagramme

\begin{tikzcd}
Z \arrow[r, "\kappa"] \arrow[dr, "\lambda\mapsto\lambda^2"'] & {Z'} \\
 & Z \arrow[u, "2\alpha"']
\end{tikzcd}

\note{sous $\kappa$, le renvoi « cf. c) ». La flèche verticale est marquée « $2\alpha$ » alors que d) nomme l'isomorphisme $\bar\alpha$ et l'oriente de $Z'$ vers $Z$ ; on transcrit tel quel.}

\marginal{cela implique que $\tilde{B}''/\tilde{U}=Z$ commutatif.}
\note{note en diagonale à gauche de b), avec un signe biffé avant $Z$.}

Nous devons maintenant exprimer que \add{dans} la situation (26) correspondant \struck{\ill{}} \add{provient} d'un (57), \struck{$S$ comme un des \ill{} $Z$} \add{\ill{}} graphe à homothéties, ayant \struck{les arêtes} \add{\ill{} d'arêtes l'ens. orienté} comme un 2-polyèdre régulier correspondant à la structure 2-polyédrale régulière de $G_0$, et comme groupe d'homothéties $Z$ ; \struck{Je \ill{} \ill{} plus \ill{} que \ill{} axiomes} \ill{} $G$ soit le groupe des automorphismes stricts des multigraphes à homothéties envisagés.
\note{les dernières lignes sont très raturées ; la lecture en est partielle.}

\page{176}
\note{toute la page est barrée de grands traits obliques ; elle est transcrite telle qu'on la lit sous les ratures.}

\struck{Donc il s'agit d'exprimer l'axiome (GH1) de (3) (p. 81), mais on y prend $f\in G$ \struck{de façon précise, \ill{}} \struck{Il est clair qu'on} doit pouvoir prendre $f\in B'$ \ill{} provenant d'un élément de $B'$ où $B'$ opère sur $S\simeq G_0/\tilde{U}$ via son opération sur $(G_0\supset\tilde{U})$. Donc cette condition}

\marginal{\struck{grâce : « conjugué » \ill{} \ill{} $G_0$ \ill{}}}

\struck{traduisant dans la terminologie de (26) : $\forall b''\in\tilde{B}''$, définissant donc une permutation $b''\tilde{U}\mapsto\ \ \tilde{U}$ de $S\simeq G_0/\tilde{U}$, par conséquent (cette op. de $G_0$ \ill{} de $G$), il existe un $b'\in B'$ (dépendant de $b''$) tel que pour tout couple $(s,t)\in S$, on ait}
\[
(58)\qquad \text{\struck{$\{s,t\}\in A_0\iff\{b'^{-1}s,\ b'^{-1}b''t\}\in A_0$}}
\]
\struck{\ill{} On \ill{} \ill{} $s=g$ ; or
$\{s,t\}\in A_0$ $\iff\exists u\in G_0$, $s=u(s_0)$, $t=u(t_0)$,
$\{b'^{-1}s,\ b'^{-1}b''t\}\in A_0$ $\iff\exists v\in G_0$, $b'^{-1}s=vs_0$, $b'^{-1}b''t=vt_0$.}

\struck{donc il faut écrire}

\struck{(58) $\forall\lambda\in Z=G_0$, $\exists f\in G$} \struck{\ill{}} \struck{tel que, $\forall s,t\in S$, on ait ($f\in B'$)}

\struck{$\{s,t\}\in A_0$} \struck{$\iff\{f^{-1}s,\ f^{-1}\lambda t\}\in A_0$}

\struck{On voit d'abord que pour $\lambda,f$ fixés,} \add{\struck{pour}} \struck{vérifier cette condition il suffit de la vérifier pour $s=s_0$ et $t$ arbitraire,} \struck{car si $s=us_0$ ($u\in G_0$),} \struck{\ill{} $t=u(t')$} \struck{on a}
\struck{$\{s,t\}\in A_0$} \struck{$\iff\{u^{-1}s,\ u^{-1}t\}\in A_0$,} \struck{$\{s_0,u^{-1}t\}$ ;}
\struck{$\{f^{-1}s,\ f^{-1}\lambda t\}\in A_0$} \struck{$\iff$} \struck{$\{u^{-1}f^{-1}uu^{-1}s,$} \struck{$u^{-1}f^{-1}uu^{-1}\lambda t\}$} \struck{$\in A_0$}
\struck{$\{u^{-1}(f^{-1})s_0,$} \struck{$u^{-1}(f^{-1})\lambda u^{-1}t\}$} \struck{$\in A_0$}

\struck{Il faut exprimer que pour $\lambda\in Z\subset G$} \struck{et $f\in G$} \struck{tels que l'image de $f$ dans} \struck{$G/G_0=Z'$} \struck{et $\alpha(\lambda)$ soient égales :} \struck{$\alpha:Z\simeq Z'$} \struck{(57 d), alors … la relation}
\note{la page s'interrompt sur ces mots.}

\page{177}\note{p. 91 de l'auteur.}
(3) de GH1 :
\[
(58)\qquad \forall s,t\in S,\ \text{on a}\ \ \{s,t\}\in A_0\iff\{f^{-1}s,\ f^{-1}\lambda t\}\in A_0 .
\]
\struck{et donc \ill{}} Cette condition ne change pas si on \struck{\ill{}} remplace $f$ par $fg$, avec $g\in G_0$, donc elle ne dépend en fait que de l'image $\dot{f}$ de $f$ dans $G/G_0\simeq Z'\ (\simeq Z)$. Quitte à changer $f$ mod $G_0$, \uncertain{SPS} $f\in B'$. Donc la variante essentielle de l'isom. $\alpha : Z\simeq\tilde{B}''/\tilde{U}\to Z'=B'/\tilde{U}$ qui « corrige » $\kappa$, c'est la validité de (58), pour $\lambda\in Z$ et $f\in B'$ reliés par $\dot{f}=\alpha(\lambda)$.
\note{le symbole devant « , avec $g\in G_0$ » est noirci ; on lit $fg$.}
\[
(59)\quad
\begin{array}{l}
\forall\lambda\in Z,\ f\in B',\ s,t\in S, \text{ tels que } \dot{f}=\alpha(\lambda) \text{ dans } Z'=B'/\tilde{U},\\
\text{on a l'équivalence}\quad \{s,t\}\in A_0\iff\{f^{-1}s,\ f^{-1}\lambda t\}\in A_0
\end{array}
\]
Nous pouvons dans cette équivalence écrire $s=us_0$, $t=vs_0$ avec $u,v\in G_0$, \struck{et les deux}. Les deux termes de l'équivalence dans (59), s'écrivant respectivement
\[
(60)\quad
\begin{cases}
\text{a) } \exists w\in G_0,\quad w(s_0,t_0)=(us_0,vs_0)\\
\text{b) } \exists w'\in G_0,\quad w'(s_0,t_0)=\bigl(f^{-1}us_0,\ \underbrace{f^{-1}\lambda vs_0}_{f^{-1}v\lambda s_0}\bigr),
\end{cases}
\]
\marginal{Ici $t_0$ est le deuxième sommet du repère d'épinglage fixé $(s_0,t_0,u_0)$}

Soit encore, en introduisant $\rho\in G_0$ tel que l'on ait
\[
t_0=\rho s_0\qquad(\rho\in G_0) \tag{61}
\]
ces termes s'écrivant aussi respectivement
\[
(62)\quad
\begin{cases}
\text{a) } \exists w\in G_0,\quad w^{-1}u\in\tilde{U},\ (w\rho)^{-1}v\in\tilde{U}\\
\text{b) } \exists w'\in G_0,\quad w'^{-1}f^{-1}u\in B',\ (w'\rho)^{-1}f^{-1}\lambda v\in B'
\end{cases}
\]
Donc la question à vérifier s'écrit
\note{la phrase reste en suspens : la page suivante ouvre un nouveau paragraphe (X), et la vérification annoncée n'est pas faite dans ce lot.}

\page{178}\note{p. 92 de l'auteur. La numérotation des formules repart à (1).}

\section{X Antipodismes généralisés, et présentation autoduale}
\note{titre souligné, de la main de l'auteur.}

1. Soit $S$ un espace homogène sous un groupe $G$. Disons que deux points de $S$ sont \add{$G$-}\emph{équivalents} s'ils ont même stabilisateur dans $G$, \struck{\ill{}} i.e. s'ils sont dans une même fibre \struck{par} \add{de} l'application
\[
\begin{array}{rcl}
S & \longrightarrow & \mathrm{Ssgr}(G)\\
s & \longmapsto & G_s=\{g\in G\mid g.s=s\}
\end{array} \tag{1}
\]
Les \add{l'une des} classes de $G$-équivalence de points \add{(correspondant à un point de $G$-orbite)} de l'espace homogène $S$ s'identifie, donc, \struck{\ill{} à l'ensemble} \add{à l'ensemble} des conjugaisons des sous-groupes de $G$ — savoir l'ensemble des $G_s$.

Soit $Z\subset\mathfrak{S}_S$ le commutant de $G$, alors $Z$ opère \emph{librement} sur $S$, et les classes de $G$-équivalence sont les $Z$-orbites dans $S$. Si $s_0\in S$, il y a donc corr. 1-1 entre l'ensemble des $s\in S$ qui sont $S$-équivalents à $s_0$, \struck{et à} \add{\uncertain{entre}} \emph{l'orbite} i.e. l'orbite de $s_0$ sous $Z$, et \add{les} \struck{\ill{}} $f\in Z$, via l'application
\[
f\longmapsto f.s_0 \tag{2}
\]
\[
Z\xrightarrow{\ \sim\ }Z.s_0=\text{ens. des pts de } S\ G\text{-équivalents à } s_0 .
\]
\note{la phrase « Les \add{l'une des} classes … » est surchargée d'ajouts ; l'ordre de lecture en est incertain. Plus bas, « $S$-équivalents » est écrit tel quel, pour $G$-équivalents.}

Notons qu'on a
\[
Z.s_0\subset S^{(G_{s_0})}\overset{\text{déf}}{=}\text{ens. des points de } S \text{ fixés par le stabilisateur de } s, \tag{3}
\]
i.e. satisfaisant $G_s\supset G_{s_0}$.
\note{dans (3), « fixés par le » est biffé et remplacé au-dessus de la ligne par un mot illisible. Le signe d'inclusion de « $G_s\supset G_{s_0}$ » est d'une lecture incertaine.}

Notons que le choix de $s_0$ comme « origine » permet d'interpréter $S$

\page{179}
comme un quotient de $G$
\[
S\simeq G/U\qquad\text{où } U=G_{s_0} \tag{4}
\]
\note{la lettre du sous-groupe est surchargée à chaque occurrence en tête de page (une autre lettre, $A$ ou $B$, réécrite) ; on lit $U$, comme dans la suite de la page.}
ce qui identifie $Z.s_0$ à l'ensemble
\[
\mathrm{Norm}_G(U)/U\subset G/U \tag{5}
\]
$Z$ lui-même étant isomorphe canoniquement au groupe quotient
\[
Z\simeq\mathrm{Norm}_G(U)/U \tag{6}
\]
l'isomorphisme étant obtenu en associant à tout $g\in G$ normalisant $U$ la « translation à droite »
\[
gU\longmapsto ga^{-1}U \tag{7}
\]
\note{dans (7), le $g$ de gauche est surchargé ; on lit « $ga^{-1}U$ » à droite.}
(pour que $ga^{-1}U$ ne dépende que de $gU\in G/U$, et ce pour un choix de $g$, il faut et il suffit que l'on ait \struck{\ill{}} \add{\uncertain{int(a)}}
\[
\mathrm{int}(a)\,U\subset U \tag{8}
\]
et pour que l'application ainsi obtenue (7) de $G/U$ dans lui-même (qui commute aux op. de $G$ et est surjective) soit bijective, il faut et il suffit que $a^{-1}$ satisfasse la même condition (8), i.e. qu'on ait
\[
\mathrm{int}(a^{-1})\,U\subset U \tag{9}
\]
qui équivaut : $U\subset\mathrm{int}(a)\,U$\note{au-dessus du signe d'inclusion de (9), un mot biffé, « the » ou approchant.}, i.e. finalement (\uncertain{vu} compte tenu de (8))
\[
\mathrm{int}(a)\,U=U\qquad\text{i.e. } a\in\mathrm{Norm}_G(U). \tag{10}
\]
\note{dans la parenthèse ouverte avant (8), « pour un choix de $g$ » surmonte un « $g=$ » ajouté ; la parenthèse ne se referme pas avant (10). Le « $g$ » et le « $a$ » de (7) ne sont pas harmonisés sur la page.}

\page{180}\note{p. 93 de l'auteur. Les lettres $U$ en exposant et dans les classes $aU$ sont, ici encore, surchargées sur une autre lettre.}
Notons que les points de $S^{G_{s_0}}=S^{U}$ s'identifiant aux points $\dot{a}=aU$ \add{de $G/U$} (dont le stabilisateur dans $G$, savoir $\mathrm{int}(a)\,U$, contient $U$, i.e. qui satisfont (9)), \struck{ils correspondent donc} \struck{à \ill{} l'unique} \add{à l'unique} application de $S$ dans lui-même
\[
f_{s_0,a}:S\longrightarrow S\quad\text{telle que}\quad
\begin{array}{l}
1)\ f_{s_0,a} \text{ commute à l'action de } G\\
2)\ f_{s_0,a}(s_0)=a,
\end{array}
\]
\struck{bijection entre l'ens.} \add{l'application $f\mapsto f(s_0)$ est une bijection}
\note{ce passage est raturé et repris par des traits de renvoi ; l'ordre de lecture retenu est incertain.}
On trouve une bijection
\[
(11)\quad
\begin{array}{l}
\text{ens. des applications de } S \text{ dans lui-même}\\
\text{qui commutent à l'action de } G
\end{array}
\xrightarrow{\ \sim\ } S^{G_{s_0}},\qquad f\longmapsto f(s_0)
\]
qui au niveau de l'espace homogène particulier \struck{identifié} identifiant
\[
S^{G_{s_0}}=S^{U}=\{aU=\dot{a}\mid a\in G,\ \mathrm{int}(a)\,U\supset U\} \tag{12}
\]
donne comme bijection l'isom. \add{$f_a$ :}
\[
aU=\dot{a}\longmapsto\bigl(\text{application } gU\mapsto ga\,U \text{ de } G/U \text{ dans lui-même}\bigr) \tag{13}
\]
\marginal{« translation à droite »}

Mais on fera attention que \struck{nous} \add{nous} avons été obligé dans (13) de \uncertain{passer} par \uncertain{inverser} la convention d'écriture (7) [où on avait l'exposant $-1$] — ce qui a comme conséquence que
\[
f_{ab}^{(d)}=f_{b}^{(d)}f_{a}^{(d)}\qquad(\text{et non } f_{a}^{(d)}f_{b}^{(d)}). \tag{14}
\]
Il y a de nombreux \uncertain{cas} où l'une des
\note{la page s'arrête sur ces mots ; la suite est hors du lot. Dans (14), l'exposant « $(d)$ » (pour « droite ») et la lettre $f$ sont d'une lecture incertaine.}

\end{document}
